\textbf{Пример 1}

Рассмотрим следующий вызов:

\t{count\_tastiness(3, [5, 2, 1])}

В этом вызове тетя должна упаковать $3$ сумки, и в кладовке есть $3$ типа печенек:
\begin{itemize}
\item $5$ печенек типа $0$, каждая из которых имеет вкусность $1$,
\item $2$ печеньки типа $1$, каджая из которых имеет вкусность $2$,
\item $1$ печенька типа $2$, которая имеет вкусность $4$.
\end{itemize}

Возможные значения $y$~--- $[0, 1, 2, 3, 4]$. Например, чтобы упаковать $3$ сумки суммарной вкусности $3$, тетя может упаковать:
\begin{itemize}
\item одну сумку содержащую три печеньки типа $0$,
\item две сумки, каждая из которых содержит одну печеньку типа $0$ и одну печеньку типа $1$.
\end{itemize}

Так как есть $5$ возможных значений $y$, функция должна вернуть $5$.

\includegraphics[scale=0.8]{1.png}
 
\textbf{Пример 2}

Рассмотрим следующий вызов:

\t{count\_tastiness(2, [2, 1, 2])}

В этом вызове тетя хочет упаковать $2$ сумки, и в кладовке есть $3$ типа печенек:
\begin{itemize}
\item $2$ печеньки типа $0$, каждая из которых имеет вкусность $1$,
\item $1$ печенька типа $1$, которая имеет вкусность $2$,
\item $2$ печеньки типа $2$, каждая из которых имеет вкусность $4$.
\end{itemize}
Возможные значения $y$~--- $[0,1,2,4,5,6]$. Так как есть $6$ возможных значений  $y$, функция должна вернуть $6$.